We describe a processing chain for spectra of the Chandrayaan-2 Large Area Soft X-ray Spectrometer (CLASS) that uses only public data. Each \SI{8}{s} spectrum is fitted on raw counts with a Poisson likelihood and a simplified fundamental-parameter model: primary fluorescence with matrix effects and incidence geometry, scattered solar X-rays, detector and filter response, and a particle background. The incident solar spectrum is approximated by an isothermal CHIANTI spectrum at the temperature measured by the GOES X-Ray Sensor (XRS) for each exposure. Because solar lines carry \LineFracSiLo--\LineFracMgHi\% of the flux that excites Mg, Al and Si, element/Si ratios depend appreciably on this spectrum. We also describe an instrumental Al K$\alpha$ component that appears when the spacecraft is sunlit from behind and whose treatment in the background model affects Al/Si. In synthetic tests Mg/Si and Al/Si are recovered with small biases, whereas Ca/Si, Ti/Si and Fe/Si are not constrained with this approach. From \NRawFiles{} exposures in \NMonths{} monthly archives, \NMapped{} spectra pass pre-declared quality criteria and give global 2\textdegree{} maps with quality layers and no infill. Al/Si correlates with the Lunar Prospector Gamma-Ray Spectrometer (GRS) on 5\textdegree{} pixels ($r=\GrsAlR$) and with nine landing-site soils ($r=\SiteAlTwoR$). Mg/Si correlates with GRS ($r=\GrsMgR$), but its landing-site comparison is inconclusive. Refitting up to \XsmNClass{} overlapping spectra from three months with solar spectra from the Chandrayaan-2 Solar X-ray Monitor changes the median Al/Si by \XsmTwoAl{} to \XsmOneAl\% and Mg/Si by \XsmTwoMg{} to \XsmOneMg\%, within the stated systematic uncertainties.
